\section{Script Reference}
\label{sec:AppendixScripts}

This appendix documents all scripts in the \filepath{scripts/} directory.

The manual has its own reproducible build helpers:
\shellcmd{bash docs/manual/build.sh} builds from checked-in vector figures;
\shellcmd{bash docs/manual/build.sh --render-figures} regenerates the Mermaid
figures before compiling. Sources and prerequisites are documented in
\filepath{docs/manual/README.md}.

\subsection{Build and install}

\begin{description}
\item[\code{build-archimesh-client.sh}]
  Convenience wrapper that calls \code{build-client-standalone.sh}.

\item[\code{build-client-standalone.sh}]
  Builds the client plugin using Tycho/Maven with the target platform.

\item[\code{build-admin-web.sh}]
  Builds the Svelte admin UI and copies output to the server's static
  resources directory.

\item[\code{install-archimesh-plugin.sh}]
  Installs the built client plugin JAR into the local Archi installation's
  dropins directory.

\item[\code{sync-archimesh-client.sh}]
  Syncs client plugin source changes without a full rebuild.

\item[\code{prepare-client-target-platform.sh}]
  Prepares the Eclipse target platform bundles required for client plugin
  compilation.

\item[\code{mint-dev-jwt.sh}]
  Mints a local dev JWT signed with the test private key.\\
  Usage: \code{./mint-dev-jwt.sh --user alice --roles admin,model\_writer}\\
  Options: \code{--expires-in <seconds>}, \code{--audience <aud>},
  \code{--print-payload}.\\
  Env overrides: \code{DEV\_JWT\_PRIVATE\_KEY}, \code{DEV\_JWT\_ISSUER},
  \code{DEV\_JWT\_EXPIRES\_IN}, \code{DEV\_JWT\_AUDIENCE}.
\end{description}

\subsection{Development infrastructure}

\begin{description}
\item[\code{dev-up.sh}]
  Starts Docker Compose infrastructure (Neo4j + Kafka) and optionally the
  server. Usage: \code{./dev-up.sh [background]}.
  Supports OIDC mode via \code{ARCHIMESH\_IDENTITY\_MODE=oidc}.

\item[\code{dev-down.sh}]
  Stops Docker Compose infrastructure.
\end{description}

\subsection{Validation and gate checks}

\begin{description}
\item[\code{run-isolated-server-gate.sh}]
  Runs all server unit, infrastructure and WebSocket tests on fresh disposable
  Neo4j/Kafka with an ephemeral HTTP test port. It requires Python 3 for run IDs
  and port selection, removes only its own Docker project, and saves logs and
  Surefire reports under \filepath{server/target/isolated-server-gate/}.
  Equivalent entry point: \shellcmd{./scripts/validate-local.sh server-isolated}.
\item[\code{run-recovery-drill.sh}]
  Packages the server and runs the opt-in \code{ProcessRecoveryIT} against
  dedicated Docker infrastructure. It force-stops only its own child server
  JVMs, waits for real writer-lease expiry, checks pending publication and
  checkpoint recovery, and interrupts an atomic v2 restore into an empty
  database. Evidence is retained under \filepath{server/target/recovery-drill/}.
  Use \code{--skip-build} to reuse an existing package. The equivalent gate
  entry point is \shellcmd{./scripts/validate-local.sh recovery}; see
  \filepath{docs/recovery-drills.md} for scenarios and cleanup.
\end{description}

\begin{description}
\item[\code{validate-local.sh}]
  Central validation entry point.\\
  Usage: \code{./validate-local.sh <mode>}\\
  Modes:
  \begin{description}
  \item[\code{fast}] Lightweight hardening and parity checks (no Docker
    infra needed). Runs \code{check-crdt-hardening.sh} and
    \code{check-notation-parity.sh}.
  \item[\code{gate}] Full client + server CRDT gate via
    \code{crdt-gate.sh}.
  \item[\code{infra}] Full gate with local Kafka/Neo4j-backed integration
    coverage. Sets \code{RUN\_LOCAL\_INFRA\_IT=true}.
  \item[\code{ws}] Full gate with websocket end-to-end coverage. Sets
    \code{RUN\_WS\_E2E\_IT=true}.
  \item[\code{recovery}] Packaged JVM crash and fresh-database restore drill
    against dedicated Docker infrastructure.
  \item[\code{server-isolated}] All server tests on fresh disposable
    Neo4j/Kafka, including WebSocket coverage; client tests run separately.
  \end{description}

\item[\code{crdt-gate.sh}]
  Runs the CRDT test suite end-to-end. First runs hardening contract checks,
  then installs the client plugin artifact, runs client CRDT tests
  (\code{OpMapperNotationInventoryTest}, \code{OpMapperConnectionBatchTest},
  \code{CrdtEntityMergeTest}, \code{CrdtPropertyMergeTest},
  \code{CrdtOrSetTest}), then runs server CRDT tests
  (\code{ArchimeshServiceTest}, optionally
  \code{LocalInfraIntegrationTest} when \code{RUN\_LOCAL\_INFRA\_IT=true}).

\item[\code{check-crdt-hardening.sh}]
  Preflight contract check that verifies presence of core CRDT test files,
  gate scripts and spec references. Fails fast when test infrastructure
  drifts.

\item[\code{check-notation-parity.sh}]
  Compares notation key inventory across: schema definitions
  (\code{schemas/ops.json}), server runtime whitelist constants
  (\code{ArchimeshService}) and client inventory test expectations
  (\code{OpMapperNotationInventoryTest}). Ensures all three agree on
  supported notation fields.

\item[\code{check-archimesh-sanity.sh}]
  Scans Archi workspace logs for critical sync issues, integrity warnings
  and connection errors. Useful for diagnosing client-side problems.
\end{description}

\subsection{Log monitoring}

\begin{description}
\item[\code{tail-archimesh-logs.sh}]
  Tails Archi workspace logs with archimesh-related filtering. Useful for
  real-time debugging during client plugin development.
\end{description}
